\originalsection{作业 5: 动量与质心}
\begin{exercise}
原题 1. 子弹 $m_1$ 以水平速率 $u$ 射入静止块 $m_2$ 并粘住. 块在动摩擦 $\mu_k$ 平面上滑 $d$ 后从高 $h$ 的桌边落下. 碰撞近似瞬时, 块尺寸可忽略, 无空气阻力. (a) 求刚能到桌边的最小初速. (b) 初速更大时求距桌边正下方的着地水平距离.
\end{exercise}
\begin{solution}
碰后 $V=m_1u/(m_1+m_2)$. 摩擦后边缘速度平方为 $V^2-2\mu_kgd$. (a) 阈值 $u_{\min}=(m_1+m_2)\sqrt{2\mu_kgd}/m_1$, 实际离边需严格超过它. (b) 下落时间 $\sqrt{2h/g}$, 故 $x_f=\sqrt{(V^2-2\mu_kgd)2h/g}$.
\end{solution}
\begin{exercise}
原题 2. 杂技演员 $m_A$ 从蹦床以速率 $v_0$ 上跳, 在高 $h_0$ 处抓住静止小丑 $m_B$, 抓取时间可忽略. 求共同达到的最大高度.
\end{exercise}
\begin{solution}
抓取前 $v_-^2=v_0^2-2gh_0$, 需该量非负且在上升段抓取. 短时竖直动量守恒给 $v_+=m_Av_-/(m_A+m_B)$. 随后上升 $v_+^2/(2g)$, 故 $h_f=h_0+[m_A/(m_A+m_B)]^2(v_0^2-2gh_0)/(2g)$.
\end{solution}
\begin{exercise}
原题 3. 胫骨压断应力约 $1.6\times10^8$ N/m$^2$, 最细面积 $3.2$ cm$^2$. 人质量 60 kg 从 2.0 m 高跳下, 着地后质心再下降 1.0 cm 而停止, 碰撞中匀减速. (a) 求碰撞时间, 两位有效数字. (b) 平均地面力. (c) 地面力与重力比, 能否在冲量近似中忽略重力? (d) 按本模型会否压断?
\end{exercise}
\begin{solution}
落地速率 $v=\sqrt{2gh_0}$, 减速度 $a=v^2/(2d)=gh_0/d$. (a) $\Delta t=2d/v=0.0032$ s. (b) $N=m(g+a)=1.18\times10^5$ N. (c) $N/(mg)=1+h_0/d=201$, 故重力冲量相对地面冲量很小. (d) 单根阈值 $FA=1.6\times10^8\times3.2\times10^{-4}=5.12\times10^4$ N; 即使两腿平均分担, 每腿约 $5.92\times10^4$ N 仍超过. 这是题设均匀减速、力分配模型的判断, 不作为个体损伤预测.
\end{solution}
\begin{exercise}
原题 4. 细杆长 $L$、质量 $M$, 沿 $x$ 轴. (a) 均匀时求相对左端的质心向量. (b) 线密度 $\lambda=\lambda_0x^3/L^3$, 求 $\lambda_0$ 与质心.
\end{exercise}
\begin{solution}
(a) $\vct R=(L/2)\uv x$. (b) $M=\lambda_0L/4$, 故 $\lambda_0=4M/L$. 质量矩 $\int_0^Lx\lambda\dd x=\lambda_0L^2/5$, 得 $\vct R=(4L/5)\uv x$.
\end{solution}
\begin{exercise}
原题 5. 粒子 $m_1,m_2$ 受距离反立方吸引. $t_0$ 时 $m_1$ 以 $v_0$ 向距它 $d$ 的静止 $m_2$ 运动, 在 $t_1$ 碰撞. 求期间粒子 1 走的距离 $L$.
\end{exercise}
\begin{solution}
取初始 $x_1=0,x_2=d$. 质心初位 $m_2d/(m_1+m_2)$, 速度 $m_1v_0/(m_1+m_2)$, 碰撞点就是质心, 因而 $L=[m_2d+m_1v_0(t_1-t_0)]/(m_1+m_2)$. 不需解奇异内力; 动量守恒已经给碰撞位置. 吸引且初始朝向对方保证粒子 1 未反向, 位移等于路程.
\end{solution}
